% !TeX root = ../main.tex
% ─────────────────────────────────────────────────────────────
%  ĐÁNH GIÁ THỰC NGHIỆM — số liệu thật từ experiments/run_eval.py
%  CSV nguồn: experiments/results/{branching,tournament,elo,timing,antirep}.csv
% ─────────────────────────────────────────────────────────────
\section{Đánh Giá Thực Nghiệm}
\label{sec:eval}

Mục này định lượng ba khía cạnh của Tactical Formation: (i) độ phức tạp
không gian trạng thái qua \emph{branching factor}, (ii) sức mạnh tương
đối của AI qua giải đấu bot-vs-bot, và (iii) hiệu năng thời gian thực.
Mục tiêu là xác nhận không gian nước đi \emph{đủ rộng và không suy biến}
(branching factor), và rằng độ sâu tìm kiếm \emph{thực sự chuyển hoá thành
sức mạnh chơi} (độ sâu lớn hơn thắng độ sâu nhỏ hơn một cách ổn định)---chính
sự phân tách theo kỹ năng này, chứ không phải branching factor, mới là bằng
chứng cho chiều sâu chiến thuật.

\subsection{Thiết Lập Thí Nghiệm}
\label{sec:eval_setup}

Mọi thí nghiệm chạy trên CPU AMD Ryzen~7 6800H (16 luồng, 7{,}4\,GB RAM),
Python~3.12 trên Linux (WSL2); tìm kiếm AI chạy đơn luồng. Mã đo đạc nằm ở
\texttt{experiments/} trong kho mã nguồn mở, mọi seed ngẫu nhiên cố định để
tái lập. Mỗi cặp đấu gồm 50 ván; để khử lợi thế tay đi đầu, hai bot luân
phiên đi đầu (mỗi bên đi đầu 25 ván). Bot dùng cùng hàm đánh giá
ở Phương trình~\eqref{eq:eval}, chỉ khác độ sâu tìm kiếm. Do mỗi nước ở độ
sâu~6 tốn vài giây trên thông dịch Python thuần (xem
Mục~\ref{sec:eval_timing}), số ván được chọn ở mức 50/cặp---đủ để chênh lệch
do độ sâu hiện ra rõ ràng mà vẫn nằm trong ngân sách tính toán.

\subsection{Branching Factor và Không Gian Trạng Thái}
\label{sec:eval_branching}

Bảng~\ref{tab:branching} báo cáo số nước đi hợp lệ trung bình (branching
factor), trên hơn $11{,}5$ nghìn thế cờ mẫu mỗi chế độ, phân theo giai
đoạn ván (đầu/giữa/cuối theo phần ba độ dài ván). Các thế cờ mẫu được sinh
bằng \emph{lối chơi ngẫu nhiên hợp lệ} (mỗi lượt chọn đều một nước trong
tập nước hợp lệ); do đó các con số phản ánh phân phối thế cờ dưới lối chơi
ngẫu nhiên, có thể lệch nhẹ so với phân phối thế cờ xuất hiện giữa hai bot
mạnh, nhưng cho ước lượng không thiên vị về bề rộng cây trò chơi.

\begin{table}[H]
  \centering
  \caption{Branching factor trung bình (số nước di chuyển hợp lệ) theo giai
  đoạn và chế độ chơi. Tính cả lựa chọn hướng kẹp, số \emph{hành động} đầy
  đủ chỉ cao hơn không đáng kể (2 người $\approx 11{,}5$; 3 người
  $\approx 14{,}2$; 4 người $\approx 12{,}1$).}
  \label{tab:branching}
  \begin{tabular}{lccc}
    \toprule
    \textbf{Giai đoạn} & \textbf{2 người} & \textbf{3 người} & \textbf{4 người} \\
    \midrule
    Đầu ván  & 11.2 & 13.7 & 11.9 \\
    Giữa ván & 11.2 & 14.1 & 11.8 \\
    Cuối ván & 11.1 & 14.1 & 11.7 \\
    \midrule
    \textbf{Trung bình ván} & 11.2 & 13.9 & 11.8 \\
    \bottomrule
  \end{tabular}
\end{table}

\noindent Branching factor trung bình ($\approx 11$--$14$) cao hơn Cờ Đam
($\approx 8$) và ngang hoặc nhỉnh hơn Reversi ($\approx 10$), và đáng chú ý
là khá \emph{ổn định} qua các giai đoạn ván---hệ quả của bất biến bảo toàn
quân số: số quân (và do đó số nước nhảy tiềm năng) không cạn dần như ở các
game có loại quân. Chế độ 3 người có branching cao nhất do bàn rộng hơn và
ba phe cùng sinh nước. Các giá trị này nằm xa dưới chặn trên giải tích
$|\mathcal A(s)|\le 1008$ (Phụ lục~\ref{app:formal}), xác nhận ràng buộc cú
nhảy thắt chặt không gian nước đi thực tế. Cần nhấn mạnh: branching factor
đo \emph{bề rộng} của cây trò chơi, không trực tiếp đo \emph{chiều sâu}
chiến thuật; nó chỉ cho thấy không gian nước đi đủ phong phú và không suy
biến về cuối ván. Bằng chứng cho chiều sâu chiến thuật---rằng các quyết
định thực sự có ý nghĩa---đến từ giải đấu ở Mục~\ref{sec:eval_tournament}.

\subsection{Giải Đấu Bot-vs-Bot}
\label{sec:eval_tournament}

Để kiểm tra AI khai thác được độ sâu tìm kiếm, chúng tôi cho các bot
khác độ sâu đấu đối kháng (Bảng~\ref{tab:tournament}).

\begin{table}[H]
  \centering
  \caption{Kết quả bot-vs-bot, chế độ 2 người (A là bên độ sâu cao hơn;
  mỗi cặp 50 ván, luân phiên đi đầu 25/25). Khoảng tin cậy 95\% của
  win-rate tính theo phương pháp Wilson. Không có ván hòa nào trong cả ba
  cặp.}
  \label{tab:tournament}
  \begin{tabular}{lcccccc}
    \toprule
    \textbf{Cặp đấu (độ sâu)} & \textbf{Số ván} & \textbf{Thắng A}
      & \textbf{Hòa} & \textbf{Thắng B} & \textbf{Win-rate A}
      & \textbf{CI 95\%} \\
    \midrule
    d=6 vs d=2 & 50 & 50 & 0 & 0 & 100.0\,\% & [92.9,\,100.0] \\
    d=6 vs d=4 & 50 & 43 & 0 & 7 & 86.0\,\% & [73.8,\,93.0] \\
    d=4 vs d=2 & 50 & 45 & 0 & 5 & 90.0\,\% & [78.6,\,95.7] \\
    \bottomrule
  \end{tabular}
\end{table}

\paragraph{Ước lượng Elo.}
Từ win-rate $p$ của bên mạnh, chênh lệch Elo ước lượng theo
$\Delta\mathrm{Elo} = -400\,\log_{10}(1/p - 1)$. Lấy bot độ sâu~2 làm
mốc 1000 Elo, ta thu được thang Elo tương đối trong Bảng~\ref{tab:elo}.
Cặp d=4 vs d=2 (90{,}0\%) cho $\Delta\mathrm{Elo}\approx{+}382$, còn cặp
d=6 vs d=4 (86{,}0\%) cho ${+}315$; ghép chuỗi, bot độ sâu~6 cao hơn độ
sâu~2 khoảng ${+}697$ Elo.\footnote{Cặp d=6 vs d=2 thắng tuyệt đối
(100\%) khiến công thức logistic phân kỳ ($p\to1$ làm
$\Delta\mathrm{Elo}\to\infty$), nên chỉ cho \emph{cận dưới} chứ không phải
một giá trị Elo hữu hạn; vì vậy chúng tôi ước lượng Elo của độ sâu~6 qua
chuỗi d=2\,$\to$\,d=4\,$\to$\,d=6 đáng tin hơn.}

\begin{table}[H]
  \centering
  \caption{Thang Elo tương đối của AI theo độ sâu (chế độ 2 người).}
  \label{tab:elo}
  \begin{tabular}{lc}
    \toprule
    \textbf{Cấu hình bot} & \textbf{Elo ước lượng} \\
    \midrule
    Độ sâu 2 (mốc) & 1000 \\
    Độ sâu 4        & 1382 \\
    Độ sâu 6        & $\approx 1697$ \\
    \bottomrule
  \end{tabular}
\end{table}

\paragraph{Lợi thế tay đi đầu.}
Để dò một khía cạnh cân bằng cơ bản của chế độ 2 người, chúng tôi cho hai
bot \emph{cùng độ sâu} ($d{=}2$) tự đánh 100 ván (seed cố định). Kết quả
đáng chú ý: bên đi \emph{sau} (Vàng) thắng 60 ván, bên đi \emph{trước}
(Xanh) thắng 36 ván, 4 hòa---tỉ lệ thắng của bên đi trước chỉ
$38{,}0\%$ (khoảng tin cậy Wilson 95\% trên các ván phân thắng
$[28{,}5,\,47{,}5]$, nằm trọn dưới $50\%$). Như vậy ở độ sâu 2, \emph{đi
sau có lợi thế} một cách có ý nghĩa thống kê---trực giác hợp lý cho một
game kẹp-lật, nơi đi sau cho phép phản đòn sau khi đối thủ đã lộ thế. Đây
là quan sát ở \emph{một} độ sâu cố định; lợi thế này có thể thay đổi theo
độ sâu và là hướng khảo sát cân bằng tiếp theo.

\subsection{Hiệu Năng Thời Gian Thực}
\label{sec:eval_timing}

Bảng~\ref{tab:timing} báo cáo thời gian suy nghĩ mỗi nước đo thực tế,
bổ sung cho các mục tiêu thiết kế ở Bảng~\ref{tab:ai_params}.

\begin{table}[H]
  \centering
  \caption{Thời gian mỗi nước đo thực tế (giây), $n=30$ nước mỗi cấu hình.}
  \label{tab:timing}
  \begin{tabular}{lccc}
    \toprule
    \textbf{Chế độ (độ sâu)} & \textbf{Trung vị} & \textbf{Bách phân vị 95}
      & \textbf{Tối đa} \\
    \midrule
    2 người (d=6) & 3.76 & 5.90 & 6.98 \\
    3 người (d=4) & 0.49 & 0.90 & 0.95 \\
    4 người (d=3) & 0.07 & 0.13 & 0.17 \\
    \bottomrule
  \end{tabular}
\end{table}

\subsection{Thảo Luận Kết Quả}
\label{sec:eval_discussion}

Ba quan sát nổi lên từ các số đo trên.

\textbf{(1) Không gian nước đi rộng và không suy biến.} Branching factor
$\approx 11$--$14$ (Bảng~\ref{tab:branching}) cho thấy mỗi lượt có nhiều
lựa chọn ý nghĩa, đồng thời \emph{ổn định} qua các giai đoạn ván---hệ quả
trực tiếp của bất biến bảo toàn quân số, khác hẳn các game có loại quân nơi
branching cạn dần về cuối ván. Trò chơi vì vậy không bị suy biến thành tàn
cuộc đơn giản. Bản thân branching không chứng minh chiều sâu chiến thuật;
nó đặt điều kiện cần (không gian không tầm thường), còn bằng chứng đủ đến
từ quan sát (2).

\textbf{(2) Độ sâu tìm kiếm chuyển hoá thành sức mạnh chơi (bằng chứng cho
chiều sâu chiến thuật).} Win-rate xếp hạng đơn điệu theo độ sâu---d=6 thắng
d=2 tuyệt đối (100\%), d=4 thắng d=2 (90\%), và d=6 thắng d=4 (86\%)
(Bảng~\ref{tab:tournament})---với \emph{không một ván hòa nào}, và các
khoảng tin cậy 95\% đều nằm trên $50\%$ một cách rõ rệt. Việc nhìn xa hơn
mang lại lợi thế hệ thống chứng tỏ các quyết định trong Tactical Formation
\emph{có hệ quả dài hạn}---dấu hiệu của chiều sâu chiến thuật---và hàm đánh
giá kết hợp tìm kiếm Minimax là ``có nghĩa'', không phải nhiễu ngẫu nhiên.
Thang Elo (Bảng~\ref{tab:elo}) định lượng khoảng cách: mỗi hai bậc độ sâu
tương ứng $\approx 300$--$380$ Elo. Cần lưu ý đây là thang \emph{nội tại}
(neo tương đối giữa các bot cùng engine), không phải Elo neo theo người
chơi hay một đối thủ ngoài; nó chứng minh sức mạnh \emph{tương đối} tăng
theo độ sâu, chứ không khẳng định chất lượng tuyệt đối của AI.

\textbf{(3) Hiệu năng thời gian thực.} Chế độ 3 và 4 người đạt dưới
$1$~giây mỗi nước (trung vị $0{,}49$\,s và $0{,}07$\,s), thoải mái cho
tương tác. Tuy nhiên chế độ 2 người ở độ sâu~6 tốn trung vị $3{,}76$\,s và
tối đa gần $7$\,s mỗi nước trên thông dịch Python thuần đơn luồng---vượt
ngưỡng phản hồi tức thì. Đây là một đánh đổi sức mạnh--tốc độ mà chúng tôi
thảo luận ở Mục~\ref{sec:limitations}, cùng các hướng giảm thiểu (hạ độ
sâu thích ứng, tìm kiếm theo thời gian, hoặc cài đặt biên dịch).

Tóm lại, thực nghiệm xác nhận cả hai mục tiêu đặt ra ở đầu mục: trò chơi
\emph{có} chiều sâu chiến thuật đo được và AI \emph{chơi hợp lý}, với chi
phí thời gian rõ ràng cho cấu hình tìm kiếm sâu nhất.
